\begin{otherlanguage}{english}
The end of Chapter XV asks how a homotopy theory can be constructed from a class of weak equivalences and a supply of small pieces. Grothendieck approaches that question in three stages. He first develops a calculus for classes recognized on bounded approximations. He then tries to turn that calculus into a general factorization theorem, discovers that the proof has lost its cardinal bound, and recovers a more structured theorem for monomorphisms. Finally, he applies the result to presheaves and returns to the axioms for a fundamental localizer and its associated derivator. The progression is mathematically instructive because the abandoned argument and the successful one use nearly the same constructions. What distinguishes them is control of the morphism left over after an attachment \citep{Grothendieck1990XV}.

The accessibility language must be read with its original conventions. A cardinal $\pi$ is an infinite bound on sizes, and a filtered indexing category “large with respect to $\pi$” has the directedness now associated with the regular cardinal $\pi^+$. At a good cardinal, an object called $\pi$-accessible is therefore $\pi^+$-presentable in current terminology. The accessibility of a \emph{class}, however, is not merely accessibility of its full subcategory. It is the embedded recognition condition $\mathrm b_\pi$: in every continuous, suitably directed presentation by objects of the chosen bounded filtration, the colimit belongs to the class if and only if cofinally many stages do. The letter $L$ adds closure under all small filtered colimits. These distinctions are essential both to the positive results of \S5 and to the defect found in \S6.

Under the standing subobject, image, and exactness assumptions, an object has a canonical presentation by its small subobjects. Proposition 5.7, its corollary, and the ensuing lemmas explain when a class can be recovered from its bounded part. The basic index is a directed poset closed under bounded directed suprema. A subset that is both cofinal and closed under those suprema can be intersected with another such subset without losing cofinality: alternately enlarge an approximation into each subset and take a countable supremum. A transfinite version gives intersections of at most $\pi$ such subsets. The same idea proves closure of L-$\pi$-accessible classes under suitably bounded intersections. Finite unions use a different fact: among finitely many subsets covering a directed set, one is cofinal. Thus finite unions and arbitrary small intersections are available for the absolute L-accessibility notion, after increasing the witness cardinal when necessary.

The flattening argument in Proposition 5.9 deserves to be read alongside these cofinality calculations. A filtered family of objects, each described by bounded pieces, is replaced by the category of all those pieces mapping into the various stages. This category is an opfibration over the original index, with filtered fibres. Its total category is filtered, and the tautological source functor has the desired colimit. Regrouping the resulting diagram by its directed subdiagrams of cardinal at most $\pi$ produces an adapted presentation. The commentary supplies the intervening finality and continuity arguments. It also records the precise limit of the written proof: the proposition advertises the weakened assumption (d$_0$), but its subobject proof uses stronger hypotheses. A complete proof is given under the full standing assumptions, and another under (d$_0$) at a good cardinal. We do not identify the remaining weakening with a disproved assertion.

The passage from objects to arrows comes from the theory of categories of sections in SGA 4, Expos\'e I, 9.22--9.24 \citep{SGA4I1972}. It requires a common cardinal bound large enough to accommodate the diagram's objects, structure maps, and equations. In the arrow category a bounded object has \emph{both} bounded source and bounded target. Theorem 6.1 attempts to show that any L-accessible class of monomorphisms containing the isomorphisms is generated, in the lifting sense, by a small subcollection. Grothendieck builds a transfinite succession of attachments inside a fixed target. At the crucial step he invokes accessibility on a continuous family whose targets remain that fixed, possibly large object. His footnotes correctly withdraw the inference: this is no longer a presentation by bounded arrows.

The distinction can be seen without homotopy theory. For a set $B$ larger than $\pi$, the arrows $S\hookrightarrow B$, indexed by the subsets $S$ of size at most $\pi$, form a continuous, highly directed diagram with colimit $\operatorname{id}_B$. None of its stages is an isomorphism, although the isomorphisms form an L-$\pi$-accessible class. This diagnoses the failed inference. To decide the proposed theorem itself requires more. A recent classification provides that decision: pure monomorphisms in the category of spans of sets are not cofibrantly generated \citep[Theorem 5.1]{CoxFeigertKamsmaMazariArmidaRosicky2026}. The development at manuscript page 84 proves that this class satisfies the manuscript's stronger recognition condition, as well as closure under pushouts, transfinite composition, and retracts. It is therefore an actual presheaf counterexample to Theorem 6.1. The cited authors prove the non-generation theorem; the commentary supplies its comparison with Grothendieck's cardinal filtration.

Section 7 succeeds by imposing structure which an arbitrary class does not have. In a presheaf category, two subobjects glued over their intersection form their union, still a subobject of the ambient object. Monomorphisms are stable under pushout and filtered colimits. Small quotient images of representables supply enough bounded pieces. If the class being constructed is $\mathrm{Mon}\cap W$, pushing out a bounded trivial monomorphism gives another weak equivalence, and cancellation in $W$ proves that the residual inclusion is again a weak equivalence. The induction can therefore continue. Well-poweredness makes it terminate, by preventing a strictly increasing chain of subobjects from having arbitrarily large length.

Once a set of generators has been obtained, the small object argument gives two weak factorization systems. The remaining recognition condition is that every map right orthogonal to all monomorphisms belongs to $W$. A retract argument then identifies these maps with the trivial fibrations, and the two factorizations fit together into a model structure. The general recognition theorem later associated with Smith, in the form proved by \citet[Theorem 1.7]{Beke2000}, places this reasoning in locally presentable categories. Its closure and solution-set hypotheses explain what survives of the question on manuscript page 85. They do not validate the unrestricted theorem of \S6.

Section 8 turns the remaining recognition condition into geometry. For a small category $A$, the category-of-elements functor $i_A:\widehat A\to\mathrm{Cat}$ transports a fundamental localizer to a class $W_A$ of presheaf maps. The subobject classifier, called the Lawvere object here, is injective with respect to monomorphisms. Its two disjoint sections provide the universal candidate for an interval. Theorem 8.4 relates universal asphericity of such an interval, triviality of maps right orthogonal to all monomorphisms, and the test-category property of every slice $A/a$. An additional asphericity condition on $A$ itself gives the global test-category conclusion. These are statements about modelling homotopy types, not automatically about preserving their cartesian products.

The mapping-cylinder proof makes the interval criterion concrete. In its presheaf application, the endpoint inclusion is a trivial monomorphism and remains so after pushout. The resulting projection is a weak equivalence by two-out-of-three. The abstract Lemma 8.7 omits that pushout condition. The commentary gives an explicit counterexample in groups, then restores the needed hypothesis. It also presents the direct interval-lifting proof, available when the morphism being tested has the right lifting property against every monomorphism; this proof validates the reduced hypothesis formulation of the interval criterion without invoking the faulty abstract lemma. The corresponding later argument is Cisinski's Lemma 1.4.13 \citep{Cisinski2006}.

Two other distinctions govern the application. First, the right properness criterion in (7.4.3) puts the fibration condition on the wrong vertical side of the pullback square; a small simplicial example detects the error. Second, the two converses in (8.2.4) belong to the manuscript's stronger excision axiom. They are not consequences of ordinary left properness and do not hold for arbitrary weak homotopy equivalences. Earlier commentary blocks indexed in DER XIIb discuss that obstruction; they are referred to here without reproducing them. Cisinski's later model theorem for accessible fundamental localizers and local test categories supplies the intended presheaf construction using the appropriate gluing results \citep[Corollary 4.2.18]{Cisinski2006}. His term “accessible localizer” means generation by a set \emph{as a localizer}; its relationship with accessibility of an arrow class is a theorem, rather than a change of spelling for L-accessibility.

The final two sections return from models to the intrinsic theory. Section 9 reorganizes the localizer axioms and studies how coproducts and connected components behave after localization. The distinction between weak Theorem A, Loc(3), and its relative version over an arbitrary base, Loc(3 bis), remains significant. The argument that a failure to preserve connected components collapses the localization to a point or to the walking arrow uses compatibility with products. The editors explicitly observe that this compatibility has not been obtained from the weak axiom alone. The commentary retains that qualification and proves the degeneration mechanism under the needed product hypothesis. Later treatments make the relative hypothesis and its consequences systematic; they should not be silently substituted for the weaker definition printed on these pages.

The decomposition problem also admits a useful direct argument. If a retract-closed class preserves connected components and a coproduct of maps belongs to it, each nonempty summand is an arrow-category retract of the whole map, by collapsing the other summands to a chosen point. An empty source summand has empty target because the total map preserves components. Thus reflection of weak equivalences by summands follows under these hypotheses. This specifies one setting in which the proposed independence of decomposition and connectivity cannot hold, while keeping the manuscript's weaker axiom system separate.

Theorem 10.1 gathers the desired conclusions for $\mathrm{HOT}_W$: homotopy Kan extensions, pointwise detection of isomorphisms, strong saturation, and invariance of locally constant diagrams. Its promised next section is absent. Modern model-category and derivator results provide the appropriate existence statements under their stated hypotheses; the commentary identifies those hypotheses and the exact published references. Locally constant diagrams require particular care, because they retain monodromy. The printed reference to Loc(5 b) in part 3(b) is stale: the prose means the fibration axiom, now numbered Loc(6 b). A rational-homology equivalence $BC_2\to e$, together with the nontrivial action of $C_2$ on a two-element set, shows why summand reflection alone cannot replace that fibration condition. The shifted statement likewise requires the nonempty indexing category indicated in the manuscript; an empty shift cannot detect a weak equivalence.

These comments distinguish source slips, missing hypotheses, valid sketches, and later theorems throughout. Results supplied by the commentators carry an explicit “ours” marker only in their development blocks; that marker records responsibility for the argument, not a claim of priority. The mathematical audit and commentary in this package are AI-produced by GPT-6 Astra (OpenAI, ChatGPT Pro mode). No subsequent assembly or editing by another system is certified here.

\end{otherlanguage}
